% GENERATED FILE. Edit canonical lessons, then run npm run generate:books.
% canonical-source-hash: fnv1a64:c2079a51487ca193
% canonical-lessons: TA-C481-totai, TA-C481-kutikal, TA-C481-nontu, TA-C481-itaru, TA-C481-tatumaru

\chapter{Legs and Stumbling}
\label{ch:ta-legs-and-stumbling}
\hlchaptermodality{481}

\begin{chapteropening}
\textbf{By the end of this chapter:} \emph{I can talk about the thighs, the heels, limping and stumbling.}

Complete the last lesson of chapter 481: I can talk about the thighs, the heels, limping and stumbling.
\end{chapteropening}

\section[\texorpdfstring{\ta{தொடை} (toṭai)}{தொடை (toṭai)}]{\ta{தொடை} (toṭai) — a thigh}

\label{lesson:TA-C481-totai}



\begin{quote}
Before the new one: say the Tamil for an eyelid, then the Tamil for to blink.
\end{quote}

\subsection*{You'll want to know: \ta{தொடை}}
\textbf{\ta{தொடை}} — \emph{toṭai} — \textquotedblleft{}a thigh\textquotedblright{}.

\begin{quote}
\ta{ஓடும்போது} \ta{என்} \ta{தொடையில்} \ta{வலி} \ta{வந்தது}.
\end{quote}

\begin{quote}
\emph{ōṭumpōtu eṉ toṭaiyil vali vantatu.}
\end{quote}

\begin{quote}
While running I got a pain in my thigh.
\end{quote}

\ta{தொடை} and \ta{தொடு} \textquotedblleft{}to touch\textquotedblright{} are different words.

\begin{cousinweb}
\textbf{\ta{தொடை}} (\emph{toṭai}) is native Dravidian (DEDR 3302); the entry also gives the form \ta{துடை} \emph{tuṭai}.
\end{cousinweb}

\begin{grammarlens}[title={What you've built}]
One more word for describing what you lived and felt.
\end{grammarlens}

\subsection*{Your turn}
\begin{itemize}
\raggedright
  \item \emph{Say it:} \emph{toṭai}
  \item \emph{Say it:} \emph{toṭai}, once more
  \item \emph{Say it:} \emph{ōṭumpōtu eṉ toṭaiyil vali vantatu.} — While running I got a pain in my thigh.
  \item \emph{Recall:} say the Tamil for an eyelid, then the Tamil for to blink, then say \emph{toṭai} again
\end{itemize}

\begin{tcolorbox}[breakable,colback=teal!4,colframe=teal!35!black,title={Before you move on}]
What does \ta{தொடை} mean? (\textquotedblleft{}A thigh\textquotedblright{}.) Say it once more.
\end{tcolorbox}
\section[\texorpdfstring{\ta{குதிகால்} (kutikāl)}{குதிகால் (kutikāl)}]{\ta{குதிகால்} (kutikāl) — the heel}

\label{lesson:TA-C481-kutikal}



\begin{quote}
Before the new one: say the Tamil for to blink, then the Tamil for a thigh.
\end{quote}

\subsection*{You'll want to know: \ta{குதிகால்}}
\textbf{\ta{குதிகால்}} — \emph{kutikāl} — \textquotedblleft{}the heel\textquotedblright{}.

\begin{quote}
\ta{புதிய} \ta{செருப்பு} \ta{என்} \ta{குதிகாலைக்} \ta{கடித்தது}.
\end{quote}

\begin{quote}
\emph{putiya ceruppu eṉ kutikālaik kaṭittatu.}
\end{quote}

\begin{quote}
My new sandals rubbed my heel.
\end{quote}

\ta{குதி} + \ta{கால்} \textquotedblleft{}foot\textquotedblright{}: in Tamil, new shoes bite (\ta{கடி}).

\begin{cousinweb}
\textbf{\ta{குதிகால்}} (\emph{kutikāl}) is native Dravidian (DEDR 1707); Telugu has \emph{gudi-kālu} for the same.
\end{cousinweb}

\begin{grammarlens}[title={What you've built}]
One more word for describing what you lived and felt.
\end{grammarlens}

\subsection*{Your turn}
\begin{itemize}
\raggedright
  \item \emph{Say it:} \emph{kutikāl}
  \item \emph{Say it:} \emph{kutikāl}, once more
  \item \emph{Say it:} \emph{putiya ceruppu eṉ kutikālaik kaṭittatu.} — My new sandals rubbed my heel.
  \item \emph{Recall:} say the Tamil for to blink, then the Tamil for a thigh, then say \emph{kutikāl} again
\end{itemize}

\begin{tcolorbox}[breakable,colback=teal!4,colframe=teal!35!black,title={Before you move on}]
What does \ta{குதிகால்} mean? (\textquotedblleft{}The heel\textquotedblright{}.) Say it once more.
\end{tcolorbox}
\section[\texorpdfstring{\ta{நொண்டு} (noṇṭu)}{நொண்டு (noṇṭu)}]{\ta{நொண்டு} (noṇṭu) — to limp}

\label{lesson:TA-C481-nontu}



\begin{quote}
Before the new one: say the Tamil for a thigh, then the Tamil for the heel.
\end{quote}

\subsection*{You'll want to know: \ta{நொண்டு}}
\textbf{\ta{நொண்டு}} — \emph{noṇṭu} — \textquotedblleft{}to limp\textquotedblright{}.

\begin{quote}
\ta{கால்} \ta{வலியால்} \ta{அவன்} \ta{நொண்டி} \ta{நொண்டி} \ta{நடந்தான்}.
\end{quote}

\begin{quote}
\emph{kāl valiyāl avaṉ noṇṭi noṇṭi naṭantāṉ.}
\end{quote}

\begin{quote}
Because of a pain in his leg he limped along.
\end{quote}

Past (I …): \ta{நொண்டினேன்}, \emph{noṇṭiṉēṉ}.

Doubled participle for a repeated action: \ta{நொண்டி} \ta{நொண்டி} \ta{நடந்தான்}.

\begin{cousinweb}
\textbf{\ta{நொண்டு}} (\emph{noṇṭu}) is native Dravidian (DEDR 3786), from lameness or tiredness.
\end{cousinweb}

\begin{grammarlens}[title={What you've built}]
One more word for describing what you lived and felt.
\end{grammarlens}

\subsection*{Your turn}
\begin{itemize}
\raggedright
  \item \emph{Say it:} \emph{noṇṭu}
  \item \emph{Say it:} \emph{noṇṭu}, once more
  \item \emph{Say it:} \emph{kāl valiyāl avaṉ noṇṭi noṇṭi naṭantāṉ.} — Because of a pain in his leg he limped along.
  \item \emph{Recall:} say the Tamil for a thigh, then the Tamil for the heel, then say \emph{noṇṭu} again
\end{itemize}

\begin{tcolorbox}[breakable,colback=teal!4,colframe=teal!35!black,title={Before you move on}]
What does \ta{நொண்டு} mean? (\textquotedblleft{}To limp\textquotedblright{}.) Say it once more.
\end{tcolorbox}
\section[\texorpdfstring{\ta{இடறு} (iṭaṟu)}{இடறு (iṭaṟu)}]{\ta{இடறு} (iṭaṟu) — to trip; to stumble}

\label{lesson:TA-C481-itaru}



\begin{quote}
Before the new one: say the Tamil for the heel, then the Tamil for to limp.
\end{quote}

\subsection*{You'll want to know: \ta{இடறு}}
\textbf{\ta{இடறு}} — \emph{iṭaṟu} — \textquotedblleft{}to trip; to stumble\textquotedblright{}.

\begin{quote}
\ta{இருட்டில்} \ta{கல்லில்} \ta{இடறி} \ta{விழுந்தேன்}.
\end{quote}

\begin{quote}
\emph{iruṭṭil kallil iṭaṟi viḻuntēṉ.}
\end{quote}

\begin{quote}
I tripped over a stone in the dark and fell.
\end{quote}

Past (I …): \ta{இடறினேன்}, \emph{iṭaṟiṉēṉ}.

What you trip over takes -\ta{இல்}: \ta{கல்லில்} \ta{இடறு}.

\begin{cousinweb}
\textbf{\ta{இடறு}} (\emph{iṭaṟu}) is native Dravidian (DEDR 437); as a noun it means an obstacle.
\end{cousinweb}

\begin{grammarlens}[title={What you've built}]
One more word for describing what you lived and felt.
\end{grammarlens}

\subsection*{Your turn}
\begin{itemize}
\raggedright
  \item \emph{Say it:} \emph{iṭaṟu}
  \item \emph{Say it:} \emph{iṭaṟu}, once more
  \item \emph{Say it:} \emph{iruṭṭil kallil iṭaṟi viḻuntēṉ.} — I tripped over a stone in the dark and fell.
  \item \emph{Recall:} say the Tamil for the heel, then the Tamil for to limp, then say \emph{iṭaṟu} again
\end{itemize}

\begin{tcolorbox}[breakable,colback=teal!4,colframe=teal!35!black,title={Before you move on}]
What does \ta{இடறு} mean? (\textquotedblleft{}To trip; to stumble\textquotedblright{}.) Say it once more.
\end{tcolorbox}
\section[\texorpdfstring{\ta{தடுமாறு} (taṭumāṟu)}{தடுமாறு (taṭumāṟu)}]{\ta{தடுமாறு} (taṭumāṟu) — to falter; to fumble}

\label{lesson:TA-C481-tatumaru}



\begin{quote}
Before the new one: say the Tamil for to limp, then the Tamil for to trip.
\end{quote}

\subsection*{You'll want to know: \ta{தடுமாறு}}
\textbf{\ta{தடுமாறு}} — \emph{taṭumāṟu} — \textquotedblleft{}to falter; to fumble\textquotedblright{}.

\begin{quote}
\ta{மேடையில்} \ta{பேசும்போது} \ta{அவள்} \ta{தடுமாறினாள்}.
\end{quote}

\begin{quote}
\emph{mēṭaiyil pēcumpōtu avaḷ taṭumāṟiṉāḷ.}
\end{quote}

\begin{quote}
She faltered while speaking on the stage.
\end{quote}

Past (I …): \ta{தடுமாறினேன்}, \emph{taṭumāṟiṉēṉ}.

Of feet, of words and of memory alike.

\begin{cousinweb}
\textbf{\ta{தடுமாறு}} (\emph{taṭumāṟu}) is native Dravidian (DEDR 3033), in an entry of words for halting and swaying.
\end{cousinweb}

\begin{grammarlens}[title={What you've built}]
One more word for describing what you lived and felt.
\end{grammarlens}

\subsection*{Your turn}
\begin{itemize}
\raggedright
  \item \emph{Say it:} \emph{taṭumāṟu}
  \item \emph{Say it:} \emph{taṭumāṟu}, once more
  \item \emph{Say it:} \emph{mēṭaiyil pēcumpōtu avaḷ taṭumāṟiṉāḷ.} — She faltered while speaking on the stage.
  \item \emph{Recall:} say the Tamil for to limp, then the Tamil for to trip, then say \emph{taṭumāṟu} again
\end{itemize}

\begin{tcolorbox}[breakable,colback=teal!4,colframe=teal!35!black,title={Before you move on}]
What does \ta{தடுமாறு} mean? (\textquotedblleft{}To falter; to fumble\textquotedblright{}.) Say it once more.
\end{tcolorbox}
